\documentclass[../../main.tex]{subfiles}
\begin{document}

{\bf[解]}

まずは回転の数について考えると，
回転軸の選び方が$3$通り，各々について右ねじと左ねじ$2$通り回転の選び方があるので，
あわせて$6$通りの回転が同様に確からしい．

\subparagraph{(1)}

最初$1$の面が上においてある状態から，$1$回目の操作で$1$が上になるのは，$1$の面を通る回転軸を選んだ場合だから
\begin{align}
  p_1 = \dfrac{1}{3}
\end{align}
である．$1$が底面にくることはないから$r_1=0$で，側面にくるのはそれ以外の場合で
\begin{align}
  q_1 = 1-p_1-r_1 = \dfrac{2}{3}
\end{align}
である．以上をまとめて
\begin{align}
  p_1=\dfrac{1}{3},\quad q_1=\dfrac{2}{3},\quad r_1=0
\end{align}
である．$\cdots$(答)

\subparagraph{(2)}

$n-1$回目の操作後から$n$回目の操作を行ったときの遷移は以下のようになる．

従って漸化式は
\begin{align}
  &
  \begin{cases}
    p_n=\dfrac{1}{3}p_{n-1}+\dfrac{1}{6}q_{n-1} \\
    q_n=\dfrac{2}{3}(p_{n-1}+q_{n-1}+r_{n-1})=\dfrac{2}{3} \\
    r_n=\dfrac{1}{6}q_{n-1}+\dfrac{1}{3}r_{n-1}
  \end{cases}
  \therefore
  &
  \begin{cases}
    p_n=\dfrac{1}{3}p_{n-1}+\dfrac{1}{9} \\
    q_n=\dfrac{2}{3} \\
    r_n=\dfrac{1}{6}q_{n-1}+\dfrac{1}{9}
  \end{cases}
\end{align}
である．$\cdots$(答)

\subparagraph{(3)}

(2)から$q_n$は定数であるから$p_n$について考えれば$r_n=1-p_n-q_n$から$r_n$は求まる．
(2)で得られた漸化式を変形して
\begin{align}
  p_{n+1}-\dfrac{1}{6}=\dfrac{1}{3}\left(p_n-\dfrac{1}{6}\right)
\end{align}
だから繰り返し用いて，初項は(1)の$p_1=\dfrac{1}{3}$と合わせて
\begin{align}
  p_n=\dfrac{1}{6}\left\{1+\left(\dfrac{1}{3}\right)^{n-1}\right\}\lim{n}{\infty}\dfrac{1}{6}
\end{align}
である．従って残る$r_n$は
\begin{align}
  r_{n+1}=\dfrac{1}{6}\left\{1-\left(\dfrac{1}{3}\right)^{n-1}\right\} \lim{n}{\infty}\dfrac{1}{6}
\end{align}
である．また，$\lim_{n\to\infty}q_n=\dfrac{2}{3}$は明白である．以上をまとめて
\begin{align}
  \lim_{n\to\infty}p_n = \lim_{n\to\infty}r_n = \dfrac{1}{6}, \quad \lim_{n\to\infty}q_n=\dfrac{2}{3}
\end{align}
である．$\cdots$(答)
\newpage
\end{document}
